📌 Relationer mellan vektorer
- Ortogonalitet: u \perp v → u \perp v
- Inte ortogonal: u \not\perp v → u \not\perp v
- Parallell: u \parallel v → u \parallel v
- Inte parallell: u \nparallel v → u \nparallel v

📌 Mängdnotation
- Tillhör: x \in A → x \in A
- Tillhör inte: x \notin A → x \notin A
- Delmängd: A \subset B → A \subset B
- Är delmängd eller lika med: A \subseteq B → A \subseteq B
- Union: A \cup B → A \cup B
- Snitt: A \cap B → A \cap B

📌 Linjär algebra‑specifikt
- Span: \text{span}\{v_1, v_2\} → \text{span}\{v_1, v_2\}
- Dimension: \dim V → \dim V
- Rang: \operatorname{rank}(A) → \operatorname{rank}(A)
- Ortogonal komplement: W^\perp → W^\perp

📌 Matriser och vektorer
- Matris:
\begin{bmatrix}
1 & 2 \\
3 & 4
\end{bmatrix}
- → \begin{bmatrix}1 & 2 \\ 3 & 4\end{bmatrix}
- Vektor:
\begin{pmatrix} x \\ y \\ z \end{pmatrix}
- → \begin{pmatrix} x \\ y \\ z \end{pmatrix}

📌 Inre produkt och norm
- Inre produkt: \langle u, v \rangle → \langle u, v \rangle
- Norm: \|v\| → \|v\|
- Absolutbelopp: |x| → |x|

📌 Tips
- Använd \operatorname{...} för att skapa egna funktionsnamn i rak stil, t.ex. \operatorname{proj}.
- I OneNote fungerar det mesta, men ibland behöver du skriva dubbla streck (t.ex. \||v\|| för norm).
- För snyggare normer och inre produkter kan du använda \left och \right för att få symbolerna att växa:
- \left\| v \right\| → \left\| v \right\|
- \left\langle u, v \right\rangle → \left\langle u, v \right\rangle

✅ Sammanfattning:
- Ortogonalitet = \perp
- Parallellitet = \parallel
- Mängdnotation = \in, \notin, \subset, \cup, \cap
- Linjär algebra = \dim, \operatorname{rank}, W^\perp, \text{span}
- Inre produkt = \langle u, v \rangle
- Norm = \|v\|